% ══════════════════════════════════════════════════════
% PART 4: VPS DEPLOYMENT (ALTERNATIVE)
% ══════════════════════════════════════════════════════

\chapter{Phase 4 --- VPS Deployment (Alternative)}

\begin{infobox}
\textbf{When to use VPS instead of Railway:}
\begin{itemize}
    \item Long-term production use (database persists permanently)
    \item Cheaper for 24/7 operation (\$3.29/month vs \$6--8/month)
    \item Full control over the server
    \item Client plans to use the bot with many users
\end{itemize}
\textbf{When to use Railway instead:}
\begin{itemize}
    \item Quick deployment for demonstration
    \item No Linux knowledge
    \item Automatic updates when code changes
\end{itemize}
\end{infobox}

\section{Step 1: Buy a VPS Server}

\subsection{Choose a Provider}

\begin{table}[H]
\centering
\begin{tabularx}{\textwidth}{|l|c|c|X|}
\hline
\textbf{Provider} & \textbf{Cost/Month} & \textbf{Link} & \textbf{Notes} \\
\hline
Hetzner & \$3.29 & hetzner.com/cloud & Cheapest; recommended \\
\hline
DigitalOcean & \$4.00 & digitalocean.com & Popular; good docs \\
\hline
Oracle Cloud & \$0.00 & cloud.oracle.com & Free forever tier; complex setup \\
\hline
\end{tabularx}
\end{table}

\subsection{Create the Server}

\begin{enumerate}[leftmargin=1.5cm]
    \item Go to your chosen provider's website
    \item Create an account and add a payment method
    \item Click \textbf{``Create Server''} or \textbf{``Create Droplet''}
    \item Select these settings:
    \begin{itemize}
        \item \textbf{OS:} Ubuntu 22.04 LTS
        \item \textbf{Location:} Singapore (closest to Bangladesh)
        \item \textbf{Size:} Smallest/cheapest plan
        \item \textbf{Authentication:} Password (easier for beginners)
    \end{itemize}
    \item Set a strong password and \textbf{write it down}
    \item Click \textbf{Create} and note the \textbf{IP address}
\end{enumerate}

\section{Step 2: Connect to the Server}

Open Windows Terminal or PowerShell:

\begin{lstlisting}[style=termstyle]
ssh root@YOUR_SERVER_IP
\end{lstlisting}

Type \texttt{yes} when asked about fingerprint, then enter your password.

\section{Step 3: Set Up the Server}

Run these commands one by one:

\begin{lstlisting}[style=termstyle, title={Update and install Python}]
sudo apt update && sudo apt upgrade -y
sudo apt install python3 python3-pip python3-venv git unzip -y
\end{lstlisting}

\section{Step 4: Upload the Code}

\textbf{From a new terminal on your PC} (keep the SSH one open):

\begin{lstlisting}[style=termstyle]
scp -r "D:\Dart Project\active_dream_bot" root@YOUR_SERVER_IP:/opt/
\end{lstlisting}

Enter the server password when asked.

\section{Step 5: Configure on the Server}

Back in your \textbf{SSH terminal} (the server):

\begin{lstlisting}[style=termstyle, title={Set up Python environment}]
cd /opt/active_dream_bot
python3 -m venv venv
source venv/bin/activate
pip install -r requirements.txt
\end{lstlisting}

\begin{lstlisting}[style=termstyle, title={Remove test data and create production config}]
rm -f .env active_dream.db
nano .env
\end{lstlisting}

Paste the client's values (same format as the \texttt{.env.example} file). Save with \textbf{Ctrl+O}, exit with \textbf{Ctrl+X}.

\section{Step 6: Test Run}

\begin{lstlisting}[style=termstyle]
source venv/bin/activate
python3 bot.py
\end{lstlisting}

Test on Telegram. If it works, press \textbf{Ctrl+C} to stop.

\section{Step 7: Set Up 24/7 Service}

\begin{lstlisting}[style=termstyle, title={Create service file}]
sudo nano /etc/systemd/system/activedream.service
\end{lstlisting}

Paste this exactly:

\begin{lstlisting}[style=envstyle, title={activedream.service}]
[Unit]
Description=Active Dream Telegram Bot
After=network.target

[Service]
Type=simple
User=root
WorkingDirectory=/opt/active_dream_bot
ExecStart=/opt/active_dream_bot/venv/bin/python3 bot.py
Restart=always
RestartSec=10
Environment=PYTHONUNBUFFERED=1
Environment=PYTHONIOENCODING=utf-8

[Install]
WantedBy=multi-user.target
\end{lstlisting}

Save and exit, then:

\begin{lstlisting}[style=termstyle, title={Start the service}]
sudo systemctl daemon-reload
sudo systemctl enable activedream
sudo systemctl start activedream
sudo systemctl status activedream
\end{lstlisting}

You should see \textbf{``active (running)''} in green.

\section{Useful Commands Reference}

\begin{table}[H]
\centering
\begin{tabularx}{\textwidth}{|l|X|}
\hline
\textbf{Command} & \textbf{What It Does} \\
\hline
\texttt{sudo systemctl status activedream} & Check if bot is running \\
\hline
\texttt{sudo systemctl restart activedream} & Restart the bot \\
\hline
\texttt{sudo systemctl stop activedream} & Stop the bot \\
\hline
\texttt{journalctl -u activedream -f} & View live logs \\
\hline
\texttt{journalctl -u activedream -n 50} & View last 50 log lines \\
\hline
\texttt{nano /opt/active\_dream\_bot/.env} & Edit configuration \\
\hline
\end{tabularx}
\end{table}

\begin{successbox}
The bot is now running 24/7 on your VPS. It will automatically restart if it crashes or if the server reboots. The database is stored permanently on the server, so no data is lost.
\end{successbox}
